% The router's plan for { products { title price reviews { rating body } } },
% drawn from the JSON the router returns under X-WG-Include-Query-Plan.
% Referenced from chapter 07. Bare tikzpicture; the figure environment, caption
% and label live at the call site.
%
% The dependency arrow runs underneath the two fetch nodes rather than between
% them: its label is wider than the gap and overprinted the boxes.
\begin{tikzpicture}[
    font=\small,
    fetch/.style={draw, rounded corners=2pt, minimum width=38mm,
                  minimum height=13mm, align=center},
    root/.style={draw, rounded corners=2pt, minimum width=28mm,
                 minimum height=8mm, align=center, fill=black!8},
    tree/.style={-{Stealth[length=2mm]}},
    dep/.style={-{Stealth[length=2mm]}, dashed},
    lbl/.style={font=\scriptsize, align=center, fill=white, inner sep=2pt},
    note/.style={font=\scriptsize\itshape, align=center}
  ]

  \node[root] (seq) at (0,0) {\code{Sequence}};

  \node[fetch] (f0) at (-3.6,-2.0)
    {\code{fetch 0}\\ \code{Single}\\ subgraph \code{catalog}};
  \node[fetch] (f1) at (3.6,-2.0)
    {\code{fetch 1}\\ \code{BatchEntity}\\ subgraph \code{reviews}};

  \draw[tree] (seq.south) -- (-3.6,-0.75) -- (f0.north);
  \draw[tree] (seq.south) -- (3.6,-0.75) -- (f1.north);

  % The dependency is what forces the order. Everything else is structure.
  \draw[dep] (-3.6,-2.65) -- (-3.6,-3.5) -- (3.6,-3.5) -- (3.6,-2.65);
  \node[lbl] at (0,-3.5)
    {\code{dependsOnFetchIds: [0]}\\
     \code{Product.id}, \code{isKey: true}};

  \node[note, anchor=north, text width=115mm] at (0,-4.1)
    {Solid lines are the plan's own structure. The dashed line is the data
     dependency underneath it: fetch 1 cannot be sent until fetch 0 has
     answered, because the key it needs is in that answer.};

\end{tikzpicture}
